\section{User Study}

We conducted a controlled user study to evaluate whether users could
understand and apply our framework to author composite visualizations. We
adopted visualization re-creation tasks because they provide consistent
targets, cover the core authoring operations, and enable a controlled
comparison between systems. We compared our system with Lyra using two VA
designs selected from the VAID corpus.

\subsection{Study Setting}
We next describe the participants, tasks, and procedure of the study.

\textbf{Tasks.}
We selected two VA designs of comparable authoring complexity from the VAID
corpus as re-creation tasks. Based on a decomposition agreed upon by
visualization experts, Task~1 contained [X] chart blocks and approximately
[M] composites, whereas Task~2 contained [Y] chart blocks and approximately
[N] composites. The composite counts are approximate because the same target
design may be reproduced using different composition strategies. Participants
were asked to reproduce the main charts and their composition while ignoring
annotations and minor stylistic details.

\textbf{Participants.}
We recruited 12 participants (P1--P12) with experience in visualization
authoring, including dashboard creation, VA system design, and visualization
research or development. They were familiar with visualization tools or
grammars such as D3, Vega, or Vega-Lite. Participants were evenly divided
into two groups of six. One group completed Task~1 with {\system} and Task~2
with Lyra, while the other used Lyra for Task~1 and {\system} for Task~2,
thereby counterbalancing the task--system assignment and system order.


\textbf{Procedure.}
Each session lasted approximately two hours. For each system, participants
received a 15-minute tutorial using a practice example distinct from the
study tasks, followed by 5 minutes of free exploration and a 25-minute
re-creation task. Before each task, we introduced the input data and explained
the visual structure of the target design. When participants were blocked,
the facilitator clarified system operations without making authoring
decisions on their behalf. During each task, we recorded task completion,
completion time, and the number of mouse clicks.
After each task, participants completed a 7-point Likert-scale questionnaire
assessing ease of learning, ease of use, authoring efficiency, support for
chart composition, and overall satisfaction. Finally, we conducted a
15-minute semi-structured interview organized around three topics:
(1) how the block granularity of {\system}, compared with that of Lyra,
affected authoring;
(2) how the composite construction processes of the two systems differed;
and (3) the reasons behind participants' questionnaire ratings. We also
collected feedback on their authoring strategies, encountered difficulties,
and suggested improvements.


\subsection{Quantitative Results}

\autoref{fig:drag-drop} summarizes the questionnaire ratings and
task performance under the two conditions. We report questionnaire ratings,
completion time, and mouse clicks using medians and interquartile ranges
because of the limited sample size. We used Wilcoxon signed-rank tests for
paired comparisons of questionnaire ratings, completion time, and mouse
clicks, and an exact McNemar test to compare task completion rates. We report
effect sizes together with the statistical significance of the differences.

\textbf{Subjective ratings.}
Participants rated {\system} higher than Lyra in [dimensions], with median
ratings of [X] and [Y], respectively. The differences in [dimensions] were
statistically significant ($p=[~]$, $r=[~]$), whereas no significant
difference was observed for [dimensions]. These results suggest that
[interpretation based on the actual results].

\textbf{Task performance.}
Participants successfully completed [X/12] tasks with {\system} and [Y/12]
with Lyra. Among the completed tasks, the median completion times were
[X] and [Y] minutes, and the median numbers of mouse clicks were [X] and
[Y], respectively. {\system} resulted in [shorter/similar] completion times
and [fewer/a similar number of] interactions than Lyra. This difference
indicates that [interpretation], although [mention nonsignificance or task
variation if applicable].

\subsection{Qualitative Results}

We analyzed participants' interview responses and organized the findings
around three themes: block-level authoring and family organization, composite
construction workflows, and the rationale underlying participants'
subjective ratings.

\textbf{Block-level authoring and family organization.}
Participants perceived different trade-offs between the block granularities
of {\system} and Lyra. [X] participants considered the blocks in {\system}
[more suitable/easier to manipulate] because [reason and representative
behavior]. In contrast, [X] participants preferred Lyra's [finer/coarser]
granularity when [scenario]. Participants also reported that the family-based
design helped them [identify related blocks, reuse configurations, or manage
structural consistency]. However, [limitations or learning difficulties, if
any]. These observations indicate that [overall interpretation].


\textbf{Composite construction workflow.}
Participants described the two systems as supporting different composition
strategies. With {\system}, they typically [describe the actual workflow],
whereas with Lyra, they [describe the corresponding Lyra workflow].
Participants noted that {\system}'s composite operations made
[relationships/layout structures] more explicit and reduced [specific
operations or difficulties]. By comparison, Lyra offered [its perceived
advantage], but required participants to [specific additional operation].
Several participants also noted that [common difficulty or limitation].
Overall, these findings suggest that [interpretation of the workflow
difference].


\textbf{Perceived usability and rating rationale.}
Participants attributed their ratings primarily to [factors]. Those who
rated {\system} highly emphasized [advantages], while lower ratings were
mainly associated with [limitations]. Ratings for Lyra were influenced by
[factors]. These explanations help contextualize the quantitative results:
[connect the interview findings to the differences or lack of differences
in questionnaire ratings].